% R152 source proposal; provenance and exact prior text remain in the proposal manifest.
\section{From a reach to a testable prediction}
Start with a movement we can measure well. For the proposed snap showcase,
the arm begins extended with the palm down. The elbow and forearm swing and
turn; the wrist and hand turn upward during the snap, then move down to rest.
The model must account for that whole movement, including the preparation,
contact and return. Finger pairing, organism-specific anatomy and timing still
need a real-reference series; a human-looking animation cannot establish that
a particular monkey can perform the same gesture.

The cup reach supplies a simpler precision comparison: relate recorded signals
to measured movement before attempting the coordinated snap. Each task still
needs its own anatomy, contact, timing and independent comparison.

That comparison connects three interacting groups: task-linked brain
observations, named-muscle observations and a musculoskeletal simulation.
A learner would relate the observations to the mechanics; independent movement
and contact measurements would test its predictions. EEG records scalp
electrical activity and surface EMG records muscle electrical activity.
A separately recorded nerve or spinal signal needs its own declared channel.
The proposed apparatus uses two arm sleeves, scalp EEG and synchronized,
calibrated left and right tripod views. Vendor, channel, bandwidth and clock
qualification remain open.

Intent to action is the target. By understanding and measuring the intermediary,
we aim to make that mapping as close to one-to-one as possible within a declared
task. Specify what counts as the intended action, the participant and session,
state and context, permitted history, tolerated error and independent reference.
The test then asks whether observing the middle reduces ambiguity and improves
prediction or qualified control. It does not assume a unique hidden intention
or claim that the best possible mapping has already been achieved.

The proposed learner must explain the recorded brain and muscle activity,
where the hand moved and when. These are proposed comparisons: no paired trial,
trained model, data acquisition, fine-tuning, reinforcement-learning rollout or
accepted checkpoint is established for this EEG/EMG/movement loop. Its sources
and permissions, physical and computational behaviour, clinical interfaces,
reference comparisons and held-out tests still need independent assessment.

The completed LINK comparison supplies a warning about what would count as
success. Its 18 runs and 720 saved epochs tested a different question; both
model-admission gates and all five benefit claims failed. Neither those runs
nor a future successful prediction would show that a trainee received sensation
or produced an action. Recording, prediction, simulated movement, sensory
delivery and physical action each need evidence of their own. EEG and EMG
recorders alone cannot write sensation, and no sensory delivery or physical
actuation is established for this proposal.

\subsection{Form, function and action need a measured middle}
A structural map constrains what a system could do. A NeuroConnectome maps
specified neural connections; its organism, coverage, resolution and source
would have to be declared. The signals passing through those connections
still have to be observed or inferred, with the difference made explicit.
To predict the next action, we need the processes that turn structural
constraints into an outcome. The proposed programme follows form to function
to action, while intent to action names the task-level translation we want
to improve.

DNA to RNA and RNA to protein put a process between a starting constraint and
its downstream product. Regulatory context, RNA processing and cell state
change what that process produces; protein structure alone does not determine
the behaviour of a whole cell. The modelling jobs differ:
\href{https://www.nature.com/articles/s41586-025-10014-0}{AlphaGenome regulatory prediction}~\cite{alphagenomeRetained},
\href{https://www.nature.com/articles/s41592-024-02201-0}{scGPT cell-state modelling}~\cite{scgptRetained} and
\href{https://www.nature.com/articles/s41586-024-07487-w}{AlphaFold 3 structure prediction}~\cite{alphafold32024}
address different parts of that chain. AlphaTranscriptome names a desired
predictive capability here. Connecting the predictors would still leave system
dynamics and the environment to account for, with experimental outcomes needed
to test the combined prediction. These modelling examples do not establish an
integrated cell programmer. The comparison uses the reported modelling roles of
AlphaGenome and scGPT; their complete methods have not been independently
assessed here.
% Source-role pointers: alphagenome2026, scgpt, alphafold3; bounded owner checks.

In the biological pathway, Brain2Spine and Spine2Body name proposed links to
test; their composition is the Brain2Body aim. The spinal cord integrates
descending and sensory signals and contributes local coordination, reflex
regulation and adaptation. This is a different mechanism from RNA processing,
but it serves the same explanatory purpose: account for the context-dependent
process between an initial constraint and its outcome. The constrained example of
\href{https://pubmed.ncbi.nlm.nih.gov/38603479/}{Spinal sensorimotor adaptation}~\cite{lavaud2024}
illustrates a spinal process; supported observations would have to test our proposed pathway. Where they are absent, the internal
states remain inferred or simulated. Physical spinal work remains read-only.
That example establishes neither unrestricted skill transfer nor validation
of our proposed pathway. The comparison draws on the abstract and
relevant author-manuscript passages; equivalence to the final published edition
and the complete methods and statistics remain unverified here.

Thought2Action relates a declared intention to an independently measured action.
Action2World follows the action to its consequence in the task environment;
Thought2World names the proposed composition. Body2World drew attention to the
consequence outside the body; Body2Action focuses on body mechanics and a
specified task outcome, while Brain2Motion focuses on the movement pathway.
These working labels identify related scopes for a test. They do not establish
interchangeable or working controllers.

The engineering route connects neural activity to musculoskeletal simulation
and then to a programmable model, toward NeuroTranslation. The proposed
movement-focused learner is a large movement model. System 3 thinking names
the proposed human--AI supervisory partnership around it. The supervisory
model would compare task strategies and diagnose where neural or mechanical
predictions miss independent measurements, then propose training or calibration
comparisons for the signal mapping and virtual body. The task learner would
perform the declared prediction or control; the supervisor would help choose
what to test. Neither role replaces the biological spinal intermediary or a
physical interface. Final held-out outcomes must remain outside rewards and
model selection.
\par\noindent Shaw and Nave's SSRN preprint,
\emph{Thinking---Fast, Slow, and Artificial}, and its
\href{https://executiveeducation.wharton.upenn.edu/thought-leadership/wharton-at-work/2026/05/thinking-fast-slow-and-artificially/}{institutional discussion}
provide conceptual context for this proposed human--AI role.~\cite{shawNaveSystem3Retained}
The comparison is limited to that framing; the complete preprint and institutional
article have not been reviewed here. Neither establishes a third anatomical brain
system or validates this programme's proposed learner or supervisor. Expert footage may supply
movement and tactical examples, with rights and camera uncertainty accounted
for; it supplies no measured expert EEG, EMG, muscle force or stimulation plan.

Action changes the outcome. Its sensory consequences inform perception and
the next intention. Independently measured movement would anchor the action
part of that loop; it would not by itself decode a unique sensation or
perception. Known stimuli, measured feedback and appropriate recordings or
reports would have to test those further mappings. The return path must
therefore be measured rather than drawn into existence. The full
Interoceptomics aim includes internal-body sensing and connectivity; external
touch and proprioceptive tasks contribute evidence but do not by themselves
solve that wider problem. A prediction, assisted execution and retained skill
without assistance answer different questions. Each wider claim must earn
its own independent comparison.

\subsection{Context connects. Grounding checks}
Within a reach, context has a concrete job: relate brain activity to muscle
activity in this task and at this time. Attention offers a candidate computation for those relationships, following
\href{https://arxiv.org/html/1706.03762v7}{Vaswani et al. (2017)}~\cite{vaswani2017} and \href{https://aclanthology.org/P19-1656/}{Tsai et al. (2019)}~\cite{tsai2019}. The model compares a representation called a query with representations called keys, then uses the resulting weights to combine their values. Masks keep unavailable observations and future
samples out. The calculation weights information; it does not establish that
attention understands all context.
% Source registry IDs: context-native-01a12f62, context-attention2017-v7, context-mult2019, context-policy4

To see what those weights mean, follow one permitted observation through the
calculation. An encoder turns a stated signal segment into a list of numbers,
called a vector. Learned calculations produce three such descriptions: the
query describes what this part of the model is looking for, a key describes
an observation it can compare with that query, and a value carries information
that observation can contribute. These are roles in a calculation. Their
entries are learned features, not automatically electrical amplitudes, muscle
forces or identified nerve cells. The original channel, time and quantity
still need their own labels.

The proposed technical method writes the operation as

\[\operatorname{Attention}(Q,K,V)=\operatorname{softmax}\!\left(QK^{\mathsf{T}}/\sqrt{d_k}+M\right)V.\]

Read it from the inside out. Q holds the query vectors, K the key vectors and
V the value vectors. A query and key have the same number of entries, their
common width $d_k$. For one query--key comparison, multiply corresponding
entries and add those products. That sum is a dot product. The representations
are learned; multiplying and adding form the numerical matching score. It is not a measured biological
connection strength or proof that the matched observation caused the movement.

The superscript T on K means transpose. It rearranges rows into columns so
the matrix multiplication can calculate a score for each query against each
key. Transposing does not reverse the recording's time order or change which
electrode supplied a value. Dividing the scores by $\sqrt{d_k}$ means dividing
by the square root of the common vector width. With comparably sized components,
a wider dot product can have a larger scale; this division counteracts that
typical growth before scores become weights. It does not calibrate the
prediction's uncertainty or establish a physical unit for a learned feature.

M is the admitted mask. For each query it specifies which observations may
contribute. In the additive form shown, an allowed entry can leave a score
unchanged, while an excluded entry can be represented by negative infinity
so that it receives zero weight. The actual numerical implementation still
needs checking. An unavailable input is not an observed zero, and a query
with no admitted keys has no admitted weighted answer; the interface must
preserve that unavailability rather than manufacture a normal-looking result.

Softmax converts the permitted scores into relative weights. It applies the
exponential function to a score, giving a positive quantity, then divides
by the sum of those quantities for that query. The permitted weights therefore
sum to one when an admitted answer exists. A larger score receives more weight;
softmax does not simply choose the largest score and discard every other
permitted observation. The final multiplication by V multiplies each value
vector by its weight and adds the results. That combined representation still
needs a declared output calculation to predict a named target in its stated
units. A neat set of weights is not yet a correct reach.

The mask must enforce the claim all the way through the calculation. For an
EEG-only prediction, withheld muscle and movement targets cannot re-enter as
queries, keys, values or features computed upstream. Future observations must
also stay out of a prediction issued before they are available. A timestamp
alone cannot establish this: filtering, alignment or scaling could already
have used a later sample. Calling the attention mask causal cannot undo that
earlier leak. Shared cues and movement artefacts still need their own controls,
even when every matrix operation is correct.

No query, key, value weights, EEG encoder or gradient route have been selected
or executed for this proposed loop. The calculation supplies an inspectable
method. Independently measured movement must still supply its test.

A pattern can fit inside the model and still describe the wrong movement.
Grounding therefore means testing the declared prediction against qualified
observations under the same conditions. Independently measured movement would
check the outcome, while muscle activity would add an observation that
constrains the brain model. Even a matching reach cannot identify a unique
private experience. The proposed paired trial would have to show which of
these observations adds useful information.

\subsection{Muscle activity tests the brain model}
Record brain and named-muscle activity during the same reach. In the proposed
paired trial, EEG would sample scalp electrical activity and surface EMG would
sample activity associated with the named muscles. Paired training could use
the muscle observation to constrain their relationship. Then withhold the
muscle target and movement reference from the EEG-only predictor: can the
permitted brain observations predict a specified muscle-activity observable
or movement?

Pairing and calibration remain proposed. Shared cues and movement artefacts
could give that predictor a shortcut, so the comparison needs controls for
both. The muscle observations would cover only the named muscles, and EMG
supplies neither afferent nerve activity nor direct force. Sensory claims
still need known stimuli and measured feedback. Several neural states could
remain compatible with even a well-matched pair; the pair could narrow a
task-specific explanation without revealing the complete hidden brain.
Checking the movement itself requires the next independent observation.
% Source registry IDs: context-native-01a12f62, context-policy4, overhaul-brantley2018
\par\noindent The candidate reference at \url{https://www.nature.com/articles/sdata2018133} has not yet been checked for exact support or reuse rights.

\subsection{Measured motion tests the virtual body}
Measure the same hand and arm from the aligned left and right tripod views.
The proposed motion-capture relay-composite would align the two point maps
in a common clock and coordinate frame, then join them into one composite
while retaining each point's origin. Before that joined map can check a
prediction, its calibration, alignment and capture precision would have to
be independently verified.

The movement map would constrain the external movement model: surface motion
and pose in the task's physical frame. To explain that motion mechanically,
the internal musculoskeletal model would describe specified bones, joints
and soft tissues. Organism, developmental stage and task would determine the
appropriate anatomical sources. X-ray projections could inform bone and
joint relationships; appropriate MRI references could inform specified soft
tissues. Registration, rights and source-specific limits would need independent
checks. No scan set is selected for this proposed comparison. Anatomy still
leaves activation, forces and neural function to test.

With a qualified composite, the musculoskeletal-simulation physics-engine's
joints, muscles, contacts and physical limits would face a measured comparison.
The simulation methods in \href{https://journals.plos.org/ploscompbiol/article?id=10.1371/journal.pcbi.1006223}{Seth et al. (2018)}~\cite{opensim2018} and
\href{https://myosuite.readthedocs.io/en/latest/}{MyoSuite documentation}~\cite{myosuiteDocsUndated} would support this comparison: fit the model on training
data, then compare its predicted reach with held-out measured movement. A force claim would require an additional force
measurement. The body's form limits possible motion; measured function would
test the mechanics and show where they need refinement.
% Source registry IDs: context-native-01a12f62, context-component-names-01a12f81, context-policy4, context-opensim2018, context-myosuite-docs, context-precision-01a12f77


A lifelike hand could expose a missed contact or a late reach. The comparison
still needs independent capture alignment, precision, mechanics and biological
fidelity; none has been demonstrated for this proposal. Captured geometry
cannot establish hidden neural state, internal forces or unique subjective
experience. It would check the declared movement prediction within those limits.

\subsection{Two pairs. One conditional learner}
The learner would have two connected jobs: relate EEG observations to muscle
activity, and relate measured movement to predicted mechanics. Large movement
model names the proposed family; deep neural network describes a possible
implementation. Large refers to intended scope, not a demonstrated parameter
count. Encoders would accept only observations permitted for the named test.
Two branches could use the attention and multimodal computations of
\href{https://arxiv.org/html/1706.03762v7}{Vaswani et al. (2017)}~\cite{vaswani2017} and \href{https://aclanthology.org/P19-1656/}{Tsai et al. (2019)}~\cite{tsai2019} to summarize those relationships
in a shared task representation, with separate output parts for each direction's targets.
% Source registry IDs: context-native-01a12f62, context-policy4, context-attention2017-v7, context-mult2019


Signal features belong to branches, and branches belong to the learner.
That hierarchy describes the computation, not two anatomical thinking systems.
A common clock, task and physical frame would connect the branches; independent
measurements would check their errors. Form constrains the prediction, and
checked function would guide refinement. Several hidden states could still
fit the same reach, so neither a shared representation nor conditional
predictions would establish an exact inverse nervous system. The architecture
has not been selected, implemented, trained or independently accepted.
To choose it, we would have to compare it with simpler baselines on
appropriately reserved observations.

\subsection{Learn the movement, then test what transfers}

A recording of a reach already contains a possible lesson for a computer: use the earlier movement to predict where the hand goes next. It need not contain EEG to pose that question. A separate proposed comparison would ask whether learning from one person's grasp-and-lift movement helps predict another person's movement. The WAY-EEG-GAL description identifies continuous recordings of lifting series as well as separate windowed trials and lift-summary tables.\citep{luciw2014way} The continuous movement could supply both the history and the later checking observation. It would not require matching AllLifts Run or trial rows, and it would create no measured neural pair between the two people. The original study of five complete weight series remains a different design.

To transfer anything useful, we must say what corresponds. Homologous structures share an evolutionary anatomical origin. Mechanically analogous structures can perform a comparable job without that shared origin: a robot gripper's opening and a human thumb--finger opening might both describe how an object is approached. That example would require a separate robotic study; the proposed two-person study would not test it. Each body keeps its own joints and possible directions of movement, positive-axis conventions, proportions, muscle--tendon roles, loads, inertia, contact and timing. Two similar-looking limbs need not generate the same forces or move at the same rate. Calculating a corresponding movement for another body is retargeting; the calculation remains derived, alongside separately identified tracking observations, estimates, predictions and simulations.

The proposed description of movement contains thirteen quantities. Three give the index sensor's position relative to the wrist, three do the same for the thumb and three for the object. One describes the distance between the two digit sensors; the final three describe wrist velocity using only information already available. Thirteen extra mask values say which features are actually available. Each person needs checked position units, sensor placement and a wrist orientation that turns positions into the intended frame without reflecting the body. After the checked conversion to metres, the middle, or median, wrist-to-index-sensor distance from that person's assigned training support supplies a positive length scale for position and velocity. This is a sensor-distance scale, not a measured bone length; development observations do not set it. Velocity is also multiplied by the same 0.1-second interval in both people. The thirteen prediction features then have no physical units. To recover velocity errors in metres per second, multiply by that person's length scale and divide by 0.1 seconds. Missing features contribute no target error; writing a masked zero does not make a missing observation measured. These features do not supply all finger joints, individual muscle forces or the person's inertia.

The model would use fifty movement frames at 100 observations per second to predict 100 milliseconds beyond the latest filtered input actually available. The source rate must first be checked as 500 observations per second. A proposed filter combining 51 present and earlier samples would precede keeping every fifth sample, producing 100 observations per second. Its linear-phase delay is 50 milliseconds. That delay stays in the account of when the information became available; moving a timestamp backward cannot make the predictor respond earlier. Filter coefficients, the frequencies it preserves, timestamp conventions, gaps and handling at recording edges still need checking. Only the named movement columns may enter. The source's kinematics matrix also contains force-plate signals, which cannot quietly become movement features.

Knowing which input was read is different from knowing whether two inputs are independent. Trustworthy published participant attribution, an unchanged exact archive entry and a checked capture can identify the participant record and recorded input without requiring one particular spelling of an ID inside the file. They do not authenticate a biological person or prove that two files came from different original acquisitions. A longer recording and its trimmed copy could have different names and file fingerprints while sharing the same movement. Shared samples, trims, derived versions and overlaps must therefore remain together as one acquisition group before filtering, cutting windows, scaling or assigning roles. The complete public non-ST roster and these relationships still need establishing. All 294 rows of the earlier AllLifts table remain unavailable for training, evaluation and assigning trials or acquisition groups. ST, competition storage and reserved LINK recordings remain closed.

Qualified P1 groups would provide the earlier movement learning. P2 groups would be divided between target training and development checking. At least two independently established P2 acquisition groups are needed. If there are $G\geq2$, a fixed hash calculation ranks their established original-parent identifiers; the first $\lfloor3G/5\rfloor$ become training groups and the rest become development groups. The floor brackets mean rounding down, so the nominal 60/40 division is explicit even when the number of groups is small. The exact identifiers, how they are written as bytes and the complete roster must be settled before ranking or seeing outcomes. The order is not chronology and need not balance lifting conditions. If uncertain ancestry forces all P2 recordings into one group, the comparison stops. Cutting that group into windows cannot create independent training and checking acquisitions. No actual group or movement segment has been assigned by this proposed identity method.

The three training arms keep the target task and model the same. B0 first learns from correctly mapped P2 training movement. B1 and B2 first learn from exactly the same P1 windows in the same order. B2 uses the proposed correct correspondence. B1 deliberately uses the thumb sensor for the index role, the object sensor for the thumb role and the index sensor for the object role, leaving the wrist anchor fixed. Both use the same correct source calibration and scale before this change, then the same correct target adapter and ordered target-training batches. Sampling gives each original acquisition group balanced treatment and records repeated as well as unique windows. The learner might undo the false correspondence; changing roles can also change which features are available. Equal numbers of updates do not give the arms equal diversity of source movements. Those limits must stay in the explanation of any difference, rather than prompting a more convenient control after results are known.

The proposed network is a single gated recurrent unit, or GRU, layer: a calculation that updates a numerical memory as it reads the sequence. It has a memory width of 32, accepts the thirteen movement features and thirteen masks as 26 inputs, and predicts thirteen outputs. That software memory is not measured biological memory. Each arm uses the same initialization for a given fixed random seed, with seeds 17, 31 and 73. Every arm/seed receives 200 earlier-learning updates and 100 target-training updates, using batches of 16 examples. Adam, the proposed weight-update rule, has learning rate 0.001, beta settings 0.9 and 0.999, epsilon $10^{-8}$ and zero weight decay. There is no clipping, changing learning-rate schedule or early stopping; the final update supplies the checkpoint. During training, squared prediction errors are averaged over observed target features; a batch with no observed targets must stop. Across three arms and three seeds this would be 2,700 updates. None has been performed for this comparison. No external checkpoint is selected, and the numerical updates, filter, precision and actual computing limits still need independent checks.

For each development acquisition group and seed, the proposed score adds the squared errors for the observed target features, divides by how many such targets were observed, then takes the square root. It then averages those group scores equally, so a long recording does not automatically outweigh a short one. A comparison with no observed targets has no valid denominator and must stop. The report must give feature counts, missing observations, errors converted back to physical units, uncertainty and a persistence forecast---continuing the latest available movement value---using the same information endpoint. Every worse group or feature remains visible. B2 would need lower normalized error than both B0 and B1 in each fixed seed, with no applicable original requirement failing, to justify another investigation. Two people and development checking do not establish performance across a population, untouched final confirmation or biological teaching.

Larger movement collections could supply more examples after their exact files, permissions, relevant movements and original-recording overlap are checked. \href{https://www.addbiomechanics.org/download_data.html}{AddBiomechanics} distinguishes recorded markers and ground reactions from fitted joint movements and torques estimated by inverse dynamics---calculating forces or torques consistent with motion and a body model. Those estimates are not new direct force measurements. The selected \href{https://amass.is.tue.mpg.de/license.html}{AMASS} asset would need its research-use, non-commercial and redistribution terms checked. The same original recording in two collections stays one source, not independent training and checking data. LaBraM's earlier training inventory flags the WAY source family, but exact checkpoint, participant and trial exposure remains unknown. A later decision to reserve a recording cannot remove its possible influence from already trained weights. No such checkpoint is selected for this comparison.

Scans have a more specific job than making the body look complete. Before adding X-ray/CT or MRI, check which anatomical sources the chosen musculoskeletal model already uses. A new scan is useful when it fills a stated gap or provides a genuinely independent check. Its person or species, scale, registration, permissions and uncertainty still need checking. Registration establishes the correspondence between that image and the body being studied. Another person's scan remains an external reference; a generic body model does not secretly measure this participant's geometry, muscle activation or forces. No scan set is selected here.

The actual eligible recordings, acquisition groups, units, frames, clocks and calibration are still unresolved for this continuous comparison. It reports no new fit or successful transfer. A movement description that proved useful could later be tested alongside genuinely paired neural and body observations, using only the inputs allowed for that prediction. The transferred or calculated target would not become a simultaneous measurement. The useful next answer is where the movement forecast improves, where it worsens and which measurement or assumption explains the difference. Neural equivalence, causal physiology, delivered sensation and a skill that remains after guidance ends would each need a further experiment.

\subsection{Pretrain. Pair. Verify. Refine}
Those comparisons start with the observations we may use. Data rights,
calibration and held-out subjects, sessions and tasks determine which can
enter the study; resolve them before processing data or choosing a model.
Open-source code does not grant blanket permission to use datasets,
checkpoints or clinical interfaces. Qualified data could first train
representations for each recording type. LaBraM, described by \href{https://arxiv.org/html/2405.18765v1}{Jiang, Zhao and Lu (2024)}~\cite{labram2024}, is one EEG encoder
candidate, not a complete brain-body writer. Synchronized paired observations could
then train the declared mapping from permitted inputs.

That first split has to survive the model's past. Before fitting a scaler,
training a tokenizer or selecting an encoder, identify the exact proposed
checkpoint and its pretraining corpus: versions, original recordings,
participants, sessions, tasks and every derivative used. The same recording
can appear in several collections; overlapping windows can appear under
different filenames. Keep those copies and derivatives with their original
recording and participant when defining the applicable splits. A new archive
name does not make the person or trial new.

The current overlap audit flags the WAY-EEG-GAL source family in LaBraM's
pretraining inventory. The exact checkpoint, participant and trial exposure
remains unknown. No LaBraM checkpoint has been selected for the separate
proposed hand/arm transfer comparison, and its target development observations
are development data throughout. A model tested there cannot be described as
independently confirmed on untouched final recordings. The audit needs the
actual source and checkpoint identities before that stronger claim can be
assessed. Unknown overlap is a limit on an independence claim; it is not a
blanket prohibition on useful pretraining or transfer.

For components fitted in this study, means, scales, tokenizer parameters,
codebooks and encoder updates must use only their assigned training support.
Development observations can make the predeclared selections; final held-out
observations must not set the preprocessing, representation, mapping, stopping
point or reward. A separately admitted external pretraining corpus can be
useful, but its exposure remains part of the model's history. Removing a record
from the later fine-tuning set cannot remove its influence from weights already
trained on it. Freezing a previously exposed model does not turn that record
back into independent final confirmation.

LaBraM's candidate sequence separates three learning problems. First, a neural
tokenizer would learn a way to describe specified signal segments. A tokenizer
is the calculation that turns those segments into model tokens. In this route,
a learned set of numerical patterns, called a codebook, supplies discrete
code identities for the segments. A code identifies a model pattern; it is not
an EEG voltage, a neuron identity, a muscle command or a word whose meaning is
already known. The codebook's learned numerical entries and the code labels
that select them are different things. Their signal support and tokenizer
version have to stay with the training record. The segment definition,
tokenizer objective and numerical implementation remain choices to qualify;
none has been selected for our paired learner.

Second, masked neural-code prediction would train an encoder using those code
targets. Hide selected signal segments from the predictor, give it the admitted
remaining context, and ask it to predict the code identities assigned to the
hidden segments by the specified tokenizer. Its output assigns scores or
probabilities to possible codes. The training loss penalizes a poor prediction
of the saved target code at the hidden positions. Computing that target from
the original training segment is legitimate supervision; leaking the hidden
segment, its target code or a feature calculated from it into the predictor
would defeat the exercise. Predicting an easy-to-leak answer would teach us
little about the representation.

Hiding a segment for learning and excluding an observation by permission serve
different purposes. A learning mask makes a stated prediction exercise; an
input-availability or causal mask limits what the named route may know.
Pretraining must declare its permitted context. An offline masked-segment
exercise does not by itself establish that an online prediction uses only
available past and present inputs. That later claim needs the entire input
route checked under its deadline, including the preprocessing.

Third, downstream fine-tuning would adapt the representation to the declared
task. Actual paired observations would supply the named EEG, muscle or movement
targets for the permitted direction. The outputs would then face their own
units, uncertainty, baselines and independent measurements. A useful code
prediction score cannot substitute for a movement error, a force measurement
or retained biological learning. Likewise, the previously reported neural
next-bin training is its own objective; it is not evidence that this staged
tokenizer route has already run.

Movement-only pretraining and unpaired representation learning remain useful
candidate routes. Transfer may use homologous structures with shared anatomical
origin or mechanically analogous structures with comparable jobs, after their
actual mechanics, scale, timing, contact and missing features are qualified.
Such training could supply a movement representation before every source has
simultaneous neural observations. Retargeted or learned movements must keep
their derived identity; they do not become measured neural--muscle--movement
pairs. The proposed transfer and the paired prediction ask complementary
questions, each with its own matched controls and independent target. Neither
is a new fit or a successful transfer reported here.

Paired prediction would give later refinement a baseline to beat. Before
optional reinforcement learning, post-training or task-specific fine-tuning,
the physical predictions and reward need checking against independent
outcomes; a reward the computer can calculate supplies no biological proof
by itself. The offline reinforcement-learning method of \href{https://arxiv.org/abs/2006.04779v3}{Kumar et al. (2020)}~\cite{cql2020} is a
candidate only where qualified actions, transitions and rewards support it: labels alone cannot tell us what an untried action
would do. Backpropagation can update differentiable neural components, but
it cannot automatically pass through every physics engine. The simulation
methods in \href{https://journals.plos.org/ploscompbiol/article?id=10.1371/journal.pcbi.1006223}{Seth et al. (2018)}~\cite{opensim2018} and \href{https://myosuite.readthedocs.io/en/latest/}{MyoSuite documentation}~\cite{myosuiteDocsUndated} do not establish
universal differentiability for this proposed route. Simpler baselines and frozen final
tests remain through every phase.
% Source registry IDs: context-native-01a12f62, context-policy4, context-labram2405, context-cql2020-abstract, context-opensim2018, context-myosuite-docs


This proposed paired training sequence has not been carried out.

\subsection{Respond quickly. Improve carefully}
A frozen accepted model would provide the proposed response: predict from
permitted past and present observations, then prepare encoded outputs for a
qualified interface. Inspecting the errors would be a separate job. That path
would calibrate and propose a candidate update, to be independently evaluated
with a restorable checkpoint before it could replace the running model.
Updating an estimate of the current state therefore differs from changing
learned weights; examining an error would not silently change the predictor.

This fast-and-slow comparison is a Kahneman-like analogy for the proposed division of work. It does not identify two biological systems or assign inherent thinking speeds to EEG and spatial data. Real-time performance would have to be measured from acquisition through filtering, alignment, inference, output, feedback and rendering, including delay tails and jitter. No accepted model, real-time demonstration or running system that adapts under these rules is established here. The loop has no measured end-to-end latency or accuracy guarantee.

That timing comparison must also count missed deadlines. An observation that
arrives too late cannot repair the response already due. Test the complete
chain while it coexists with the declared healthy workloads, and report the
conditions as well as delay tails and jitter. Timing one model calculation
alone cannot establish that acquisition, prediction, output and feedback meet
the task's deadline together. These remain proposed measurements, not a
real-time result from this loop.
% Source registry IDs: context-native-01a12f62, context-policy4

\subsection{Write sensation. Enact action}
What reaches the trainee when the computer has made its prediction? A sensory
transducer would have to deliver a sensory output; a controller and actuator
would have to produce action in a virtual, robotic or later qualified biological
body. Input and output describe the computer boundary. Afferent and efferent
describe physiological direction relative to the central nervous system.

A display or haptic interface supplies sensory feedback, which differs from direct neural writing. Direct cortical writing could have a sensory role without being a peripheral afferent signal, just as a robot command is not an anatomical efferent nerve. EEG and EMG recorders do not write, and predicting an EEG-like pattern does not deliver sensation. Each proposed output needs an independently qualified interface, permission for its use and a measurement of what it delivers. No such physical interface, stimulation, acquisition or actuation is established by this proposal.

Comparing a recorder with a sensory or action device requires evidence about what each measures or delivers. The clinical source could not be read, so it supplies no methods evidence here. Implanted precedents cannot establish that EEG or EMG recorders are equivalent to devices that deliver sensation or produce an action. Those physical additions remain unpriced and outside the acquisition and virtual-modelling scope.
% Source registry IDs: context-write-native-01a12f66, context-native-01a12f62, context-policy4

\subsection{Close the loop with evidence}
Return to the cup with two independent checks: withheld body observations
for the brain-to-body prediction, and the recorded neural target for the
conditional estimate from the body. Agreement between the models or a high
reward could still accompany a wrong answer. Neither establishes a unique
inverse brain or someone's private experience. A later sensory or action
interface would need a further measurement of what it delivered, under which
conditions, with what error and delay.

Approximation gets us towards plausibility and virtualization gets us to
precision. Independent measurements, repeatability, causal checks and stated
error bounds must earn that precision; the proposal supplies neither perfect
biology nor an automatic result. The original held-out and real-reference
controls would test where each prediction works and where it fails. The
retained-learning controls ask a further question: what remains when useful
guidance stops, with any necessary movement bridge kept constant?

Physical sensory writers, transducers and actuators remain unpriced additions
outside the acquisition and virtual-modelling costs. Those costs cannot tell us
the price of delivering sensation or action. No completed loop, training or
physical-writing result is claimed here. \emph{Follow the action, then test the
prediction} develops the tasks that would challenge the mechanism;
\emph{Measure both ends. Test the translation.} gives the detailed measurement
and validation protocol. Comparable observations, predictions and errors would make those tests possible.
% Source registry IDs: context-native-01a12f62, context-write-native-01a12f66, context-precision-01a12f77, context-policy4, context-spec15, context-glossary13, b2b-circular-checklist-original
